%==============================================================================
%  STAT 5200 - Applied Econometrics I
%  Homework 3
%
%  Compile with:  pdflatex "HW 3 - Report.tex"  (run twice for refs)
%
%  Directory layout assumed by this file:
%     HW 3 - Report.tex
%     R/        (scripts, for Q3 and Q4)
%     figures/  (generated figures, for Q3)
%     output/   (captured console output, for Q3 and Q4)
%==============================================================================
\documentclass[11pt]{article}

\usepackage[margin=1in]{geometry}
\usepackage[T1]{fontenc}
\usepackage[utf8]{inputenc}
\usepackage{amsmath,amssymb}
\usepackage{booktabs}
\usepackage{graphicx}
\usepackage{xcolor}
\usepackage{listings}
\usepackage[font=small,labelfont=bf]{caption}
\usepackage[hidelinks]{hyperref}

%------------------------------------------------------------------------------
%  Code presentation
%
%  Code is set as selectable text on a fixed character grid, source and console
%  output are given visually distinct styles, and source lines are kept to 76
%  characters so that nothing wraps.
%------------------------------------------------------------------------------
\definecolor{codekeyword}{RGB}{0,86,148}
\definecolor{codestring} {RGB}{163,21,21}
\definecolor{codecomment}{RGB}{106,115,125}
\definecolor{codebg}     {RGB}{248,248,250}
\definecolor{coderule}   {RGB}{220,220,226}

\lstdefinestyle{Rcode}{
  language        = R,
  basicstyle      = \ttfamily\footnotesize,
  keywordstyle    = \color{codekeyword}\bfseries,
  stringstyle     = \color{codestring},
  commentstyle    = \color{codecomment}\itshape,
  numbers         = left,
  numberstyle     = \ttfamily\scriptsize\color{codecomment},
  numbersep       = 9pt,
  backgroundcolor = \color{codebg},
  frame           = single,
  rulecolor       = \color{coderule},
  framesep        = 4pt,
  xleftmargin     = 15pt,
  xrightmargin    = 0pt,
  columns         = fixed,
  basewidth       = 0.5em,
  keepspaces      = true,
  tabsize         = 2,
  showstringspaces = false,
  breaklines      = true,
  breakatwhitespace = true,
  postbreak       = \mbox{\textcolor{codekeyword}{$\hookrightarrow$}\space},
  captionpos      = b,
  abovecaptionskip = 6pt,
}

\lstdefinestyle{Routput}{
  language        = {},
  basicstyle      = \ttfamily\footnotesize,
  keywordstyle    = \ttfamily,
  commentstyle    = \ttfamily,
  stringstyle     = \ttfamily,
  numbers         = none,
  backgroundcolor = \color{white},
  frame           = leftline,
  rulecolor       = \color{coderule},
  framesep        = 6pt,
  xleftmargin     = 20pt,
  breaklines      = true,
  captionpos      = b,
  abovecaptionskip = 6pt,
}

\lstset{style=Rcode}
\renewcommand{\lstlistingname}{Listing}

% Short inline code fragments.
\newcommand{\code}[1]{\texttt{#1}}

% Operators used throughout.
\DeclareMathOperator{\Var}{Var}
\DeclareMathOperator{\Cov}{Cov}
\DeclareMathOperator{\Avar}{Avar}
\DeclareMathOperator{\sd}{sd}
\DeclareMathOperator*{\plim}{plim}

% Placeholder marker. Every occurrence should be gone before submission;
% search the source for "\todo" to find what is still outstanding.
\newcommand{\todo}[1]{\textcolor{red}{\textbf{[TODO:} #1\textbf{]}}}

%------------------------------------------------------------------------------
\title{\vspace{-2em}Homework 3\\[0.3em]
       \large STAT 5200 --- Applied Econometrics I --- Fall 2026}

\author{Apostolos Kyratzis - 13130833}
\date{October 15, 2026}

\begin{document}
\maketitle

%==============================================================================
\section{Q1: Reflection}
\label{sec:q1}
%==============================================================================

\paragraph{Are my answers in line with the solutions?}
\todo{read \code{Homework Solutions/Homework 2/HW 2 - Solutions.pdf} and state
whether the answers agree; list any question missed}

\paragraph{What could I have done better?}
\todo{points for improvement}

%==============================================================================
\section{Q2: Asymptotic properties of the OLS estimator}
\label{sec:q2}
%==============================================================================

Throughout, the model is
\[
y_i = x_i\beta + u_i, \qquad i = 1,\dots,N,
\]
with $y_i$ a scalar and $x_i$ a $1 \times k$ row vector, and the estimator is
\[
\hat\beta = \left(\sum_{i=1}^{N} x_i'x_i\right)^{-1}
            \left(\sum_{i=1}^{N} x_i'y_i\right).
\]
We assume OLS1 and OLS2 and that $(y_i,x_i)$ is i.i.d.\ over $i$. Each part
below flags where the assumptions are used.

\paragraph{(1) Consistency of $\hat\beta$.}
\todo{show $\plim \hat\beta = \beta$; Notes 03, pp.~4--5}

\paragraph{(2) The asymptotic distribution of $\hat\beta$.}
\todo{derive $\sqrt{N}(\hat\beta - \beta) \xrightarrow{d} N(0, A^{-1}BA^{-1})$;
Notes 03, p.~6 and Appendix 1}

\paragraph{(3) A consistent estimator of the asymptotic variance.}
\todo{the sandwich estimator; Notes 03, Appendix 2, eq.~(40)}

\paragraph{(4) The standard error of $\hat\beta$.}
\todo{square root of the $k$-th diagonal entry of $\hat{V}/N$; Notes 03, pp.~6--7}

%==============================================================================
\section{Q3: Returns to education in the CPS85 data}
\label{sec:q3}
%==============================================================================

The data set is \code{CPS85.txt} from Lab 2. It is part of the course material
distributed by the instructor and is therefore not tracked in this repository;
the script reads it from its original location by relative path.

\paragraph{(1) Estimating the wage equation.}
The model is
\[
\log(\text{wage}_i) = \beta_0 + \beta_1 \text{educ}_i + \beta_2 \text{exper}_i
 + \beta_3 \text{exper}_i^2 + u_i .
\]
\todo{estimate by OLS and report the coefficient table}

\paragraph{(2) Normal quantile plot of the residuals.}
\todo{produce the Q-Q plot, then judge whether $u_i$ looks normal and say
whether non-normality matters for the later analysis}

\paragraph{(3) The parameter $\theta$.}
Increasing education by one year while experience falls by one year gives
\[
\theta = 100\left[\exp\!\left(\beta_1 - \beta_2 + (1-2\,\text{exper})\beta_3\right) - 1\right].
\]
\todo{comment on the construction of $\theta$}

\paragraph{(4) An estimator for $\theta$, and its consistency.}
\todo{plug-in estimator $\hat\theta$; argue consistency}

\paragraph{(5) $\hat\theta$ across experience levels.}
\todo{compute $\hat\theta$ for exper $= 0,5,\dots,55$ and discuss}

\paragraph{(6) How standard errors and confidence intervals would be computed.}
\todo{discussion only, no computation required}

\paragraph{(7) Standard errors and confidence intervals (optional).}
\todo{optional, not graded; delete this paragraph if not attempted}

%==============================================================================
\section{Q4: Monte Carlo simulation and the empirical coverage rate}
\label{sec:q4}
%==============================================================================

The population model is
\[
y = \beta_0 + \beta_1 x + \beta_2 x^2 + u,
\qquad
\beta_0 = 1,\;\; \beta_1 = 1,\;\; \beta_2 = 0.2,
\]
with $u \sim t_7$ and $x \sim N(1,4)$, and $N = 100$.

\paragraph{(1) Sampling one data set.}
\todo{draw $Y_N^{(1)}$ and $X_N^{(1)}$}

\paragraph{(2) The OLS estimate.}
\todo{compute $\hat\beta = [\hat\beta_0, \hat\beta_1, \hat\beta_2]'$}

\paragraph{(3) A 95\% confidence interval for $\beta_1$.}
\todo{build the interval from the asymptotic distribution of Q2(2), and define
the indicator $I^{(m)}$}

\paragraph{(4) The empirical coverage rate over $M = 10{,}000$ replications.}
\todo{report the coverage rate and say whether it matches expectations, by
reference to the assumptions}

\paragraph{(5) The heteroskedastic model.}
With the error scaled by $|x|/2$,
\[
y = \beta_0 + \beta_1 x + \beta_2 x^2 + \frac{|x|}{2}\,u ,
\]
\todo{report the coverage rate, compare it with part (4), and explain by
reference to the assumptions}

%==============================================================================
\section{Q5: Paper question}
\label{sec:q5}
%==============================================================================

\paragraph{(1) The paper.}
\todo{a paper from one of the five journals listed in Homework 1, containing
"OLS" or "ordinary least squares" and reporting OLS results; confirm the
publication venue before writing the citation}

\paragraph{(2) Citation.}
\todo{author, title, year, journal, volume, issue, pages, DOI}

\paragraph{(3) The main research question.}
\todo{one paragraph}

\paragraph{(4) The empirical model in our notation.}
Writing their specification as $y_i = x_i\beta + u_i$:
\todo{identify $y_i$, identify $x_i$ in the main specification, give the number
of variables in $x_i$, and give the total number of observations}

\paragraph{(5) Do the authors state the assumptions behind OLS?}
\todo{report what the paper does and does not state}

\paragraph{(6) A plausible story under which $E[x_i u_i] = 0$ fails.}
\todo{one paragraph, no more than five sentences, as critical as possible}

%==============================================================================
\appendix
\section{Complete R scripts}
%==============================================================================

\subsection{\texorpdfstring{\code{R/q3\_cps85\_regression.R}}{R/q3\_cps85\_regression.R}}
\lstinputlisting[style=Rcode,caption={\code{R/q3\_cps85\_regression.R}
--- Q3, the CPS85 wage equation.},label={lst:app-q3}]%
{R/q3_cps85_regression.R}

\subsection{\texorpdfstring{\code{R/q4\_coverage\_rate.R}}{R/q4\_coverage\_rate.R}}
\lstinputlisting[style=Rcode,caption={\code{R/q4\_coverage\_rate.R}
--- Q4, Monte Carlo simulation and the empirical coverage rate.},
label={lst:app-q4}]%
{R/q4_coverage_rate.R}

\end{document}
